This section takes a small station, a bus with two feeder breakers, from an empty
drawing to an exported calculation report. Each step refers to the part of the User
Guide (Section~\ref{sec:userguide}) that covers it in detail.

\subsection*{Step 1 --- Start a layout}
Click \menu{File > New} (\keys{Ctrl + N}) and type a name in the \ui{Project} field at
the top right. If you have the one-line as a PDF, click \menu{File > Import PDF} and
choose the file. Its first page appears under the grid, at drawing scale
(Section~\ref{sec:files}).

\subsection*{Step 2 --- Create the equipment}
The Equipment Library starts empty. Click \menu{Edit > Create Component} and create:

\begin{itemize}
  \item a \textbf{Bus}, named \emph{BUS-1}, with its kV rating;
  \item a \textbf{Circuit Breaker}, named \emph{CB-F}, with its trip coil, close coil
        and motor currents, for example Trip Coil~1 = \SI{6}{A}, Trip Coil~2 =
        \SI{5}{A}, Close Coil = \SI{10}{A};
  \item a \textbf{Current Transformer}, named \emph{CT-F}.
\end{itemize}

Each one appears in the Equipment Library (Figure~\ref{fig:create}).

\subsection*{Step 3 --- Place and connect}
Drag \emph{BUS-1} from the library onto the drawing, then drag in two breakers and two
CTs, one of each per feeder. Place each feeder's CT between the bus and its breaker, so
that the CTs form the boundary of the bus zone.

To draw a wire, press the mouse button on a terminal (one of the dark dots on a
symbol) and drag to another terminal. To connect to the bus, release the wire
anywhere on the bus bar; a connection node is added where it lands
(Section~\ref{sec:connecting}).

\subsection*{Step 4 --- Check the data}
Click each breaker. The \ui{Selected Element Properties} panel shows the currents it
was created with; changes made here apply to that breaker only. Scroll down to the
breaker settings and choose the DC supply for each coil, for example Trip Coil~1 on
\emph{DC 1} and Trip Coil~2 on \emph{DC 2} (Section~\ref{sec:properties}).

\subsection*{Step 5 --- Validate}
Click \menu{Analysis > Validate} and fix any open terminals or shorts it reports. End
each feeder that leaves the drawing with an \textbf{Incoming AC Line} symbol, so that
its last terminal is connected (Section~\ref{sec:validate}).

\subsection*{Step 6 --- Evaluate and configure the zone}
Click \menu{Analysis > Run Evaluation}. The tool proposes a \emph{BUS-1 Bus Zone}
bounded by the two CTs and outlines it on the drawing. The zone is marked
\emph{review required} because no operating loads have been assigned yet.

Click \menu{Analysis > Configure Zones} (Figure~\ref{fig:zonedialog}). For the bus
zone:
\begin{enumerate}
  \item make sure both feeder CTs are checked as boundary CTs;
  \item click \ui{Add Load} once for each coil that operates. For example, add each
        breaker's Trip Coil~1 and Trip Coil~2 in the \emph{Primary trip} stage, and a
        close coil in the \emph{Reclose} stage if the scheme recloses;
  \item check \ui{I have verified this zone's CT boundaries and operating actions},
        then click \ui{Save}.
\end{enumerate}
The evaluation runs again automatically (Section~\ref{sec:zones}).

\subsection*{Step 7 --- Review and export}
The Analysis Output now reads \textbf{COMPLETE}. It lists each stage's total current
and the peak on each DC supply, and the zone outline becomes solid. Click
\menu{File > Export Results} and choose \ui{Save Calculation PDF} to save the report
(Section~\ref{sec:reports}). Save the layout with \keys{Ctrl + S}.

\begin{note}
To try the tool without building a drawing, open one of the provided layouts
(Section~\ref{sec:providedlayouts}) and click \ui{Run Evaluation}.

\end{note}
